% GENERATED FILE. Edit canonical lessons, then run npm run generate:books.
% canonical-source-hash: fnv1a64:afa544f5318e365b
% canonical-lessons: AR-C436-mirab, AR-C436-kushk, AR-C436-shubbak, AR-C436-hammal, AR-C436-tanaqqala

\chapter{Around the station}
\label{ch:ar-around-the-station}
\hlchaptermodality{436}

\begin{chapteropening}
\textbf{By the end of this chapter:} \emph{I can say a garage, a kiosk, a ticket window, a porter and to travel about.}

Complete the last lesson of chapter 436: I can say a garage, a kiosk, a ticket window, a porter and to travel about.
\end{chapteropening}

\section[\texorpdfstring{\ar{مرآب} (mirʾāb)}{مرآب (mirʾāb)}]{\ar{مرآب} (mirʾāb) — a garage; a car park}

\label{lesson:AR-C436-mirab}



\begin{quote}
Before the new one: say the Arabic for a pedestrian crossing, then the Arabic for to turn around.
\end{quote}

\subsection*{You'll want to know: \ar{مرآب}}
\textbf{\ar{مرآب}} — \emph{mirʾāb} — \textquotedblleft{}a garage; a car park\textquotedblright{}.

\begin{quote}
\ar{ترك} \ar{أبي} \ar{السيارة} \ar{في} \ar{المرآب}.
\end{quote}

\begin{quote}
\emph{taraka abī as-sayyāra fī al-mirʾāb.}
\end{quote}

\begin{quote}
My father left the car in the garage.
\end{quote}

\begin{cousinweb}
\textbf{\ar{مرآب}} (\emph{mirʾāb}) is a modern coinage, usually traced to the native root \emph{r-ʾ-b}, \textquotedblleft{}to repair, to mend\textquotedblright{} (Lane), in the tool-and-place pattern \emph{mifʿāl}: a place where cars are kept and seen to.
\end{cousinweb}

\begin{grammarlens}[title={What you've built}]
One more word for getting around.
\end{grammarlens}

\subsection*{Your turn}
\begin{itemize}
\raggedright
  \item \emph{Say it:} \emph{mirʾāb}
  \item \emph{Say it:} \emph{mirʾāb}, once more
  \item \emph{Say it:} \emph{taraka abī as-sayyāra fī al-mirʾāb.} — My father left the car in the garage.
  \item \emph{Recall:} say the Arabic for a pedestrian crossing, then the Arabic for to turn around, then say \emph{mirʾāb} again
\end{itemize}

\begin{tcolorbox}[breakable,colback=teal!4,colframe=teal!35!black,title={Before you move on}]
What does \ar{مرآب} mean? (\textquotedblleft{}A garage; a car park\textquotedblright{}.) Say it once more.
\end{tcolorbox}
\section[\texorpdfstring{\ar{كشك} (kushk)}{كشك (kushk)}]{\ar{كشك} (kushk) — a kiosk}

\label{lesson:AR-C436-kushk}



\begin{quote}
Before the new one: say the Arabic for to turn around, then the Arabic for a garage.
\end{quote}

\subsection*{You'll want to know: \ar{كشك}}
\textbf{\ar{كشك}} — \emph{kushk} — \textquotedblleft{}a kiosk\textquotedblright{}.

\begin{quote}
\ar{اشتريت} \ar{جريدة} \ar{من} \ar{الكشك}.
\end{quote}

\begin{quote}
\emph{ishtaraytu jarīda min al-kushk.}
\end{quote}

\begin{quote}
I bought a newspaper at the kiosk.
\end{quote}

\begin{cousinweb}
\textbf{\ar{كشك}} (\emph{kushk}) is usually traced to Persian \emph{kushk}, \textquotedblleft{}a pavilion\textquotedblright{}; English \emph{kiosk} comes from the same word.
\end{cousinweb}

\begin{grammarlens}[title={What you've built}]
One more word for getting around.
\end{grammarlens}

\subsection*{Your turn}
\begin{itemize}
\raggedright
  \item \emph{Say it:} \emph{kushk}
  \item \emph{Say it:} \emph{kushk}, once more
  \item \emph{Say it:} \emph{ishtaraytu jarīda min al-kushk.} — I bought a newspaper at the kiosk.
  \item \emph{Recall:} say the Arabic for to turn around, then the Arabic for a garage, then say \emph{kushk} again
\end{itemize}

\begin{tcolorbox}[breakable,colback=teal!4,colframe=teal!35!black,title={Before you move on}]
What does \ar{كشك} mean? (\textquotedblleft{}A kiosk\textquotedblright{}.) Say it once more.
\end{tcolorbox}
\section[\texorpdfstring{\ar{شباك} (shubbāk)}{شباك (shubbāk)}]{\ar{شباك} (shubbāk) — a ticket window; a window}

\label{lesson:AR-C436-shubbak}



\begin{quote}
Before the new one: say the Arabic for a garage, then the Arabic for a kiosk.
\end{quote}

\subsection*{You'll want to know: \ar{شباك}}
\textbf{\ar{شباك}} — \emph{shubbāk} — \textquotedblleft{}a ticket window; a window\textquotedblright{}.

\begin{quote}
\ar{اشترى} \ar{أبي} \ar{التذاكر} \ar{من} \ar{الشباك}.
\end{quote}

\begin{quote}
\emph{ishtarā abī at-tadhākir min ash-shubbāk.}
\end{quote}

\begin{quote}
My father bought the tickets at the ticket window.
\end{quote}

Plural: \ar{شبابيك}, \emph{shabābīk}.

\begin{cousinweb}
\textbf{\ar{شباك}} (\emph{shubbāk}) is built on the native Arabic root \emph{sh-b-k}, \textquotedblleft{}to interlace, to entwine\textquotedblright{} (Lane). A \emph{shubbāk} was first a latticed window.
\end{cousinweb}

\begin{grammarlens}[title={What you've built}]
One more word for getting around.
\end{grammarlens}

\subsection*{Your turn}
\begin{itemize}
\raggedright
  \item \emph{Say it:} \emph{shubbāk}
  \item \emph{Say it:} \emph{shubbāk}, once more
  \item \emph{Say it:} \emph{ishtarā abī at-tadhākir min ash-shubbāk.} — My father bought the tickets at the ticket window.
  \item \emph{Recall:} say the Arabic for a garage, then the Arabic for a kiosk, then say \emph{shubbāk} again
\end{itemize}

\begin{tcolorbox}[breakable,colback=teal!4,colframe=teal!35!black,title={Before you move on}]
What does \ar{شباك} mean? (\textquotedblleft{}A ticket window; a window\textquotedblright{}.) Say it once more.
\end{tcolorbox}
\section[\texorpdfstring{\ar{حمال} (ḥammāl)}{حمال (ḥammāl)}]{\ar{حمال} (ḥammāl) — a porter}

\label{lesson:AR-C436-hammal}



\begin{quote}
Before the new one: say the Arabic for a kiosk, then the Arabic for a ticket window.
\end{quote}

\subsection*{You'll want to know: \ar{حمال}}
\textbf{\ar{حمال}} — \emph{ḥammāl} — \textquotedblleft{}a porter\textquotedblright{}.

\begin{quote}
\ar{حمل} \ar{الحمال} \ar{حقيبتي} \ar{الثقيلة}.
\end{quote}

\begin{quote}
\emph{ḥamala al-ḥammāl ḥaqībatī ath-thaqīla.}
\end{quote}

\begin{quote}
The porter carried my heavy bag.
\end{quote}

The pattern \emph{faʿʿāl} names a trade, as in \emph{najjār}, \textquotedblleft{}a carpenter\textquotedblright{}.

\begin{cousinweb}
\textbf{\ar{حمال}} (\emph{ḥammāl}) is built on the native Arabic root \emph{ḥ-m-l}, \textquotedblleft{}to carry\textquotedblright{} (Lane).
\end{cousinweb}

\begin{grammarlens}[title={What you've built}]
One more word for getting around.
\end{grammarlens}

\subsection*{Your turn}
\begin{itemize}
\raggedright
  \item \emph{Say it:} \emph{ḥammāl}
  \item \emph{Say it:} \emph{ḥammāl}, once more
  \item \emph{Say it:} \emph{ḥamala al-ḥammāl ḥaqībatī ath-thaqīla.} — The porter carried my heavy bag.
  \item \emph{Recall:} say the Arabic for a kiosk, then the Arabic for a ticket window, then say \emph{ḥammāl} again
\end{itemize}

\begin{tcolorbox}[breakable,colback=teal!4,colframe=teal!35!black,title={Before you move on}]
What does \ar{حمال} mean? (\textquotedblleft{}A porter\textquotedblright{}.) Say it once more.
\end{tcolorbox}
\section[\texorpdfstring{\ar{تنقل} (tanaqqala)}{تنقل (tanaqqala)}]{\ar{تنقل} (tanaqqala) — to travel about; to move from place to place}

\label{lesson:AR-C436-tanaqqala}



\begin{quote}
Before the new one: say the Arabic for a ticket window, then the Arabic for a porter.
\end{quote}

\subsection*{You'll want to know: \ar{تنقل}}
\textbf{\ar{تنقل}} — \emph{tanaqqala} — \textquotedblleft{}to travel about; to move from place to place\textquotedblright{}.

\begin{quote}
\ar{تنقلنا} \ar{بين} \ar{المدن} \ar{بالقطار}.
\end{quote}

\begin{quote}
\emph{tanaqqalnā bayna al-mudun bi-l-qiṭār.}
\end{quote}

\begin{quote}
We travelled between the cities by train.
\end{quote}

Present: \ar{يتنقل}, \emph{yatanaqqal}.

\begin{cousinweb}
\textbf{\ar{تنقل}} (\emph{tanaqqala}) is built on the native Arabic root \emph{n-q-l}, \textquotedblleft{}to move from place to place\textquotedblright{} (Lane). It is Form V, \textquotedblleft{}to keep moving on\textquotedblright{}.
\end{cousinweb}

\begin{grammarlens}[title={What you've built}]
One more word for getting around.
\end{grammarlens}

\subsection*{Your turn}
\begin{itemize}
\raggedright
  \item \emph{Say it:} \emph{tanaqqala}
  \item \emph{Say it:} \emph{tanaqqala}, once more
  \item \emph{Say it:} \emph{tanaqqalnā bayna al-mudun bi-l-qiṭār.} — We travelled between the cities by train.
  \item \emph{Recall:} say the Arabic for a ticket window, then the Arabic for a porter, then say \emph{tanaqqala} again
\end{itemize}

\begin{tcolorbox}[breakable,colback=teal!4,colframe=teal!35!black,title={Before you move on}]
What does \ar{تنقل} mean? (\textquotedblleft{}To travel about; to move from place to place\textquotedblright{}.) Say it once more.
\end{tcolorbox}
